\section{Exécutif national}
\label{executif-national}

\subsection{Composition}
\label{executif-national-composition}

\begin{enumerate}
 \item L'exécutif national sera composé de :
  \begin{enumerate}
   \item La présidence
   \item La vice-présidence académique (VPA)
   \item La vice-présidence aux communications (VPC)
   \item La vice-présidence externe (VPE)
   \item La vice-présidence aux finances (VPF)
   \item La vice-présidence aux services (VPS)
  \end{enumerate}
\end{enumerate}

\subsection{Élection}
\label{executif-national-election}

\begin{enumerate}
 \item L'exécutif national sera élu lors de l'assemblée générale annuelle.
\end{enumerate}

\subsection{Éligibilité}
\label{executif-national-eligibilite}

\begin{enumerate}
 \item La personne candidate doit faire partie d'une association membre au moment de sa candidature.
\end{enumerate}

\subsection{Mandat}
\label{executif-national-mandat}

\begin{enumerate}
 \item Le mandat de l'exécutif national commencera à la date de son élection et expirera à la clôture de l'assemblée générale annuelle de l'année suivante.
 \item À la fin de son mandat, la personne à la vice-présidence aux finances fera la transition vers le rôle de vice-présidence aux finances sortante, avec un mandat se terminant le 31 août.
\end{enumerate}
